\chapter{Explicit solitons and translating graphs}\label{sec:examples-solitons}
We work with smooth immersed hypersurfaces $\Sigma^n\subset\mathbb R^{n+1}$, $n\geq1$, and the convention $\vec H=-H\nu$, $H=\operatorname{div}_\Sigma\nu$. A parametrized family is a mean curvature flow when its normal velocity is $\vec H$; tangential velocity only changes parametrization. Definitions \ref{def:self-shrinker} and \ref{def:mcf-other-solitons} fix the scale-one shrinker and expander equations. For translators we compare with \cite{HIMW19}, Sections 3--4, converting the source's downward translation to upward translation when writing graphs.

\begin{proposition}[The scale and time interval of a homothetic soliton]\label{prop:example-homothetic-clock}
Suppose an immersion $X:\Sigma\to\mathbb R^{n+1}$ satisfies $\vec H=cX^\perp$ for a constant $c\in\mathbb R$. Then
\[
 X_t=\sqrt{1+2ct}\,X,\qquad 1+2ct>0,
\]
is a mean curvature flow in the normal-velocity sense. For $c=-1/2$ it shrinks toward time $t=1$; for $c=1/2$ it expands from time $t=-1$; for $c=0$ it is stationary. These are statements about smooth times only.
\end{proposition}
\begin{proof}
Under a dilation by $r>0$, the tangent and normal spaces keep their directions and the mean curvature vector becomes $r^{-1}\vec H$. The normal velocity of $r(t)X$ is $r'X^\perp$. Thus $rr'=c$ suffices. The solution with $r(0)=1$ is the displayed positive square root. This verifies the scaling and the maximal interval on which that formula is positive; it does not prescribe a weak continuation through the zero-scale time. Compare Proposition \ref{prop:mcf-parabolic-rescaling}.
\end{proof}

\begin{proposition}[Products preserve a soliton equation]\label{prop:example-soliton-product}
Let $\Sigma^m\subset\mathbb R^{m+1}$ be smooth. For $k\geq0$, the product $\Sigma\times\mathbb R^k\subset\mathbb R^{m+k+1}$ has mean curvature vector $(\vec H_\Sigma,0)$. It preserves a homothetic equation $\vec H=cx^\perp$. If $\Sigma$ translates with velocity $v$, the product translates with every velocity $(v,w)$, $w\in\mathbb R^k$; the extra component is tangent.
\end{proposition}
\begin{proof}
Use an orthonormal tangent frame on $\Sigma$ followed by the $k$ Euclidean coordinate vectors. The second fundamental form in the new directions and mixed directions is zero. Its trace is therefore unchanged. The normal space has no component in the product factor, giving $(x,y)^\perp=(x^\perp,0)$ and $(v,w)^\perp=(v^\perp,0)$. This proves all the assertions. In particular a translation vector is not uniquely determined by normal motion when a Euclidean product factor is present.
\end{proof}

\begin{example}[Planes, spheres, and cylinders have different motion laws]\label{ex:soliton-model-comparison}
An affine plane is a translator with velocity $v$ exactly when $v$ is tangent to it, in which case its image is stationary. It is a scale-one shrinker or expander centered at zero exactly when it passes through zero. The round sphere of radius $\sqrt{2m}$ in $\mathbb R^{m+1}$, $m\geq1$, satisfies the shrinker equation, and its product with $\mathbb R^k$ gives a shrinking cylinder. No round sphere centered at zero satisfies the scale-one expander equation.
\end{example}
\begin{proof}
The mean curvature of a plane is zero. Its normal position component is its constant displacement from the parallel plane through zero, proving the plane claims. On a radius-$R$ sphere, $\vec H=-(m/R)\nu$ and $x^\perp=R\nu$. The shrinker equation is $m/R=R/2$, whereas the expander equation would require $-m/R=R/2$, impossible for positive $R$. Apply Proposition \ref{prop:example-soliton-product} for cylinders. The evolving sphere radius is computed in Example \ref{ex:shrinking-sphere}; the stationary plane's apparent tangential translation must not be mistaken for nonzero normal speed.
\end{proof}

\begin{proposition}[The upward translator equation for a graph]\label{prop:translator-graph-example}
Let $u\in C^2(\Omega)$, $\Omega\subset\mathbb R^n$ open, and $c\in\mathbb R$. Its graph translates with velocity $ce_{n+1}$ if and only if
\[
 \operatorname{div}\frac{Du}{\sqrt{1+|Du|^2}}
 =\frac{c}{\sqrt{1+|Du|^2}},
 \qquad
 \left(\delta_{ij}-\frac{u_i u_j}{1+|Du|^2}\right)u_{ij}=c.
\]
With upward unit normal $\nu=(-Du,1)/\sqrt{1+|Du|^2}$, the scalar curvature is $H=-c/\sqrt{1+|Du|^2}$.
\end{proposition}
\begin{proof}
The upward graph normal has tangential divergence $H=-\operatorname{div}(Du/\sqrt{1+|Du|^2})$. The normal component of the velocity is $(ce_{n+1})^\perp=(c/\sqrt{1+|Du|^2})\nu$. Equating this with $-H\nu$ gives the divergence equation. Expanding the divergence and multiplying by $\sqrt{1+|Du|^2}$ gives the second form. Equivalently substitute $u(x,t)=u(x)+ct$ in Proposition \ref{prop:graphical-mcf}. The negative sign of $H$ here follows from the chosen upward normal.
\end{proof}

\begin{proposition}[All translating graphs of curves with nonzero vertical speed]\label{prop:grim-reaper-speed-family}
Let $c>0$. Every connected graph $y=u(x)$ translating upward with speed $c$ is a restriction of
\[
 u(x)=-\frac1c\log\cos(c(x-x_0))+b,
 \qquad |x-x_0|<\frac{\pi}{2c},
\]
for constants $x_0,b$. The full graph is properly embedded and complete; a restriction to a proper open subinterval is incomplete at any new finite endpoint. Its curvature magnitude is $c\cos(c(x-x_0))$, and its total curvature is $\pi$.
\end{proposition}
\begin{proof}
Proposition \ref{prop:translator-graph-example} reduces to $u''=c(1+(u')^2)$. Writing $p=u'$, integration gives $\arctan p=c(x-x_0)$ after adjusting $x_0$, and a second integration yields the formula. This also identifies the maximal interval between consecutive vertical asymptotes. The arc length element is $ds=\sec(c(x-x_0))\,dx$; its integral diverges at both endpoints. Height also tends to infinity there, giving properness. At an added interior endpoint both slope and arc length density remain bounded, giving incompleteness. Finally $|\kappa|=u''/(1+(u')^2)^{3/2}=c\cos(c(x-x_0))$, so $|\kappa|ds=c\,dx$ integrates to $\pi$. The speed-one case retains its earlier entry at Example \ref{ex:mcf-grim-reaper}; \cite{HIMW19}, Example 1.6, uses the reflected downward version.
\end{proof}

\begin{example}[Tilted grim-reaper cylinders]\label{ex:tilted-grim-reaper}
For $a\in\mathbb R$, put $A=1+a^2$ and
\[
 u(x,y)=-A\log\cos(x/\sqrt A)+ay,
 \qquad |x|<\frac\pi2\sqrt A,\quad y\in\mathbb R.
\]
This is a proper complete embedded translator with velocity $e_3$. Its strip width is $\pi\sqrt{1+a^2}$. The direction $(0,1,a)$ is tangent everywhere and leaves the graph invariant under translation. Products give higher-dimensional examples.
\end{example}
\begin{proof}
Writing $\xi=x/\sqrt A$, we have $u_x=\sqrt A\tan\xi$, $u_y=a$, $u_{xx}=\sec^2\xi$, and the other second derivatives vanish. Since $1+|Du|^2=A\sec^2\xi$,
\[
 \left(1-\frac{u_x^2}{1+|Du|^2}\right)u_{xx}
 =\frac{1+a^2}{A\sec^2\xi}\sec^2\xi=1.
\]
Proposition \ref{prop:translator-graph-example} proves the equation. A bounded ambient image bounds $x,y,u$; as $x$ tends to a strip endpoint with bounded $y$, the logarithmic term diverges. Thus inverse images of compact sets stay in a compact subdomain, proving properness. A proper smooth immersion into Euclidean space is complete: any intrinsic Cauchy sequence has bounded image and a convergent subsequence in its compact inverse image, and smooth local coordinates give intrinsic convergence. Graphs are embedded. The tangent-direction and product assertions follow directly. This is the upward form of \cite{HIMW19}, equation (3.2), with $a=\tan\theta$.
\end{proof}

\begin{proposition}[A translator's height is a strict elliptic test]\label{prop:translator-height-example}
Let $X:\Sigma^n\to\mathbb R^{n+1}$ satisfy $\vec H=v^\perp$, and put $h=\langle X,v\rangle$. Pointwise,
\[
 \nabla_\Sigma h=v^T,\qquad \Delta_\Sigma h=|v^\perp|^2,
 \qquad \Delta_\Sigma e^h=|v|^2 e^h.
\]
There is no nonempty compact boundaryless translating hypersurface in Euclidean space, including the case of zero translation velocity.
\end{proposition}
\begin{proof}
Differentiating $h$ tangentially gives its gradient. The immersion identity $\Delta_\Sigma X=\vec H$ gives $\Delta h=\langle\vec H,v\rangle=|v^\perp|^2$. The exponential chain rule adds $|\nabla h|^2=|v^T|^2$, yielding the final identity. If $v\ne0$, its integral on a compact boundaryless hypersurface contradicts the divergence theorem. If $v=0$, the immersion is minimal and $\Delta_\Sigma|X|^2=2n$, again impossible after integration. Neither completeness nor properness is required for the pointwise identities.
\end{proof}

\begin{proposition}[Translation as a weighted area problem]\label{prop:translator-weighted-area-example}
A smooth immersion satisfies $\vec H=v^\perp$ if and only if it is stationary for the weighted area $\int_\Sigma e^{\langle x,v\rangle}\,dA$ under all compactly supported normal variations. Stationarity is understood locally even when the total weighted area is infinite.
\end{proposition}
\begin{proof}
For normal initial velocity $Y$, the ordinary area derivative is $-\langle\vec H,Y\rangle dA$ and the derivative of the weight is $e^{\langle x,v\rangle}\langle v,Y\rangle$. Therefore
\[
 \delta\mathcal A_v(Y)=\int_\Sigma
 e^{\langle x,v\rangle}\langle v^\perp-\vec H,Y\rangle\,dA.
\]
Compact support makes this expression finite and removes boundary terms. It vanishes for all normal tests precisely when the vector in parentheses is zero. In \cite{HIMW19}, Proposition 3.2, the velocity is $-e_{n+1}$ and the weight is correspondingly $e^{-x_{n+1}}$; a Gaussian weight would instead describe a shrinking soliton.
\end{proof}

\begin{proposition}[The radial equation at the tip of a bowl translator]\label{prop:radial-translator-tip}
Let $n\geq2$, and suppose $u(x)=f(|x|)$ is a smooth graph near zero translating with velocity $e_{n+1}$. For $r>0$ its radial equation is
\[
 \frac{f''}{1+(f')^2}+\frac{n-1}{r}f'=1.
\]
At the center, $f'(0)=0$ and $f''(0)=1/n$. If $f$ is $C^4$, then
\[
 f(r)=f(0)+\frac{r^2}{2n}+\frac{r^4}{4n^3(n+2)}+o(r^4).
\]
These are necessary local identities; they do not by themselves prove that the solution extends to a complete entire bowl.
\end{proposition}
\begin{proof}
The Hessian of a radial function has eigenvalues $f''$ in the radial direction and $f'/r$ in the $n-1$ perpendicular directions. The gradient is radial, so substitution in Proposition \ref{prop:translator-graph-example} gives the ODE. Smooth radiality gives $Du(0)=0$ and $D^2u(0)=f''(0)I$; the graph equation at zero is then $nf''(0)=1$. Write $f=f(0)+Ar^2+Br^4+o(r^4)$ with the differentiated remainders furnished by $C^4$ regularity and radial symmetry. The constant equation gives $2nA=1$. The $r^2$ coefficient is $4(n+2)B-8A^3=0$, proving the formula. The complete rotational bowl is the example discussed in \cite{HIMW19}, Section 4; its global construction requires a separate ODE argument.
\end{proof}

\begin{example}[A paraboloid is not an exact bowl translator]\label{ex:paraboloid-not-translator}
For $n\geq2$, the function $u(x)=|x|^2/(2(n-1))$ does not solve the unit-speed translator equation, even though its growth has the quadratic scale associated with the bowl model. In fact its translator operator equals
\[
 \left(\delta_{ij}-\frac{u_i u_j}{1+|Du|^2}\right)u_{ij}
 =1+\frac{1}{(n-1)(1+r^2/(n-1)^2)}>1.
\]
\end{example}
\begin{proof}
For this radial graph, $f'=r/(n-1)$ and $f''=1/(n-1)$. Substitute into the operator in Proposition \ref{prop:radial-translator-tip}. The tangential term is one and the radial term is the additional positive fraction. In particular its Hessian at zero is $I/(n-1)$, whereas a smooth exact bowl tip must have Hessian $I/n$. An asymptotic shape must not be substituted for an exact solution.
\end{proof}
